\section{Expediente NIST y significado del puerto para cada arquitectura}
\label{sec:puertos-nist}
Suponemos una única fuente física y un presupuesto de adquisiciones por
consumo. Para usar crédito NIST después de Toeplitz, el puerto posterior
a Toeplitz debe estar incluido en la evaluación 90B con un crédito de
entropía explícito. Una fuente con acondicionador Toeplitz no vetted
puede recibir crédito parcial. En la frontera aquí elegida, con Toeplitz
como único acondicionador, ese crédito no alcanza la calificación NIST de
\emph{full entropy}. Para los consumos que la requieren añadimos
AES-CMAC vetted conforme a §3.2.2.2. Para otra frontera de fuente validada que ya entregue
full entropy en su puerto, 90C admite usar esa salida directamente
\cite[§3.1.5]{nist90b}\cite{nist-vetted}\cite[§§3.1--3.2, 6.3(1)--(7)]{nist90c}.

\paragraph{AES/SHA.}
El núcleo de §§996--998 de AIS combina primitivas normalizadas, pero
su tupla no corresponde a Hash\_DRBG, HMAC\_DRBG ni CTR\_DRBG. La exigencia
de 90C §6.3(2) se refiere al \emph{mecanismo DRBG} aprobado y validado.
Por tanto, las identidades ITS y la calificación algorítmica del ejemplo
bajo AIS no proporcionan una vía 90C RBG3 para ese mismo puerto
AES/SHA. La fuente y el acondicionador pueden servir además a otra
construcción NIST con asignación independiente de adquisiciones.
Ésta tendría su propio puerto y expediente \cite[§§996--998]{ais}\cite[§6.3]{nist90c}.

\paragraph{CTR directo.}
El mecanismo 90A sin DF admite el perfil de un bloque con
$D=0^{256}\|A$. La siembra y la resiembra tienen una interfaz distinta de
$D$ y requieren 384 bits de material con full entropy en el perfil elegido.
Una valoración bajo 90C de la construcción RBG2(P) atribuye salida
seudoaleatoria a la fuerza instanciada si se cumplen sus demás requisitos.
El método 1 de recuento admite sólo la entropía física validada para la
siembra y resiembra. El material se consume una vez. Si el DRBG forma
parte de un RBG3 y se accede directamente a él, §6.3(8) aplica las reglas
RBG2(P) a ese acceso.
Añadir CMAC después de Toeplitz permite analizar la coordenada final de
su entrada y presupuestar full entropy para los consumos que la requieren.
La salida directa de Generate continúa siendo distinta de los procesos
RBG3(XOR) y RBG3(RS) de 90C §§6.4 y 6.5
\cite[tabla 3 y §10.2.1]{nist90a}\cite[§§5.3(1)--(3),(6)--(7),(9),(11),6.2--6.5]{nist90c}.

\subsection{RBG3(RS): conservación en la última coordenada de la resiembra}
En esta variante se fija AES-256, sin DF, $s=256$, entrada adicional
ausente tanto para reseed como para Generate y petición externa $n=128$.
Sea $S$ el estado anterior a la resiembra y $W=H\|L$ su material de
384 bits con $|H|=256$, $|L|=128$. SP 800-90A Update evalúa con el
estado anterior tres bloques AES consecutivos. Llamemos $g_K(S)$ a
los dos primeros concatenados y $g_V(S)$ al tercero. Se obtiene
\[
 k'=g_K(S)\oplus H,\qquad v'=g_V(S)\oplus L .
\]
SP 800-90C §6.5.1.2 solicita al DRBG \emph{256 bits}, avanza por dos
bloques y ejecuta el Update posterior. Después devuelve el prefijo
de 128 bits. Por ello
\[
 Y=\AES_{k'}(\operatorname{inc}(v')),\qquad
 L=g_V(S)\oplus\operatorname{inc}^{-1}(\AES_{k'}^{-1}(Y)).
\]
Para $S,H$ fijos, el mapa de $L$ a $Y$ es una permutación aun al dar
la vuelta el contador. Si $Z=(S,H,E)$, el lema de cambio de variable da
\[
 \hm(Y\mid Z)=\hm(L\mid Z),\qquad
 \dist(P_{YZ},\ud_{128}\otimes P_Z)
 =\dist(P_{LZ},\ud_{128}\otimes P_Z).
\]
Esta premisa de conservación se refiere sólo a los últimos 128 bits
condicionados al prefijo y al estado. La especificación NIST sigue
pidiendo material de siembra/resiembra de 384 bits con full entropy en
el perfil sin DF. Si $H$ es fijo y $L$ uniforme, los dos bloques
internos juntos tienen soporte a lo sumo $2^{128}$, incluso cuando el
primero es uniforme. No se extiende la identidad al par generado ni
al estado posterior divulgado \cite[§§6.3(5), 6.5.1.2, ap. A.2, tabla 4]{nist90c}.

\paragraph{Efecto en un expediente dual sobre un puerto.}
El mecanismo RS hace influir la aportación física en el estado mediante la
resiembra, lo que aborda la condición estructural de AIS §381. El Update
posterior parte del contador \emph{tras el segundo bloque}. La tupla
directa demostrada para un bloque termina tras el primero. Una petición
externa de 128 bits no elimina esta diferencia. Para evaluar también el
puerto RS bajo AIS se deben especificar la nueva tupla y la granularidad
de .7, justificar su análisis autónomo para el tamaño de petición que
realmente admite, y delimitar qué números se califican como internos.
Con la interfaz de Generate de 90C, el DRBG produce y devuelve al
envoltorio dos números internos completos antes de truncar. Para este
diseño contamos conservadoramente 256 bits internos de CTR por resiembra
y 128 bits externos, con balance nominal $384/256$. AIS distingue el
número interno listo para salida del subconjunto efectivamente entregado
(§§73, 106, 128--131 y glosario). Si el solicitante propone otro límite
de TOE en el que el segundo bloque nunca está disponible como salida del
RBG completo, debe especificar esa función de salida y el avance del
estado que realmente se ejecuta. La compresión efectiva (§388), la cota de PTG.3.5 para esos
números y el propio expediente de fuente requieren justificación.
La conservación demostrada aquí se aplica al primer bloque. En una
implementación con acceso compartido al estado, 90C §6.5.2 exige completar
cada petición RS antes de la siguiente. Si la última operación fue RS,
una petición directa al DRBG que comparta ese estado exige resiembra
previa con al menos $s+64$ bits de entropía y ejecución indivisible
frente a otras peticiones. El perfil RS estudiado usa su puerto propio
\cite[§§73, 106, 128--131, 385--388 y glosario]{ais}\cite[§6.5.2(1)--(3)]{nist90c}.

\subsection{RBG3(XOR): influencia de la rama física y coste}
En el flujo normativo, $F$ procede del puerto físico full entropy
y $B$ del DRBG. Dados $B,E$, la transformación $F\mapsto Y=F\oplus B$
conserva exactamente su distancia a uniforme y su minentropía. Las
adquisiciones de $F$ son distintas de las de siembra y resiembra.
En la conexión ordinaria, con entradas auxiliares fijas, sustituir
$F$ por $F'$ conserva el estado posterior del DRBG: el XOR actúa sobre
la salida que recibe la aplicación. Por ello, identificar $F$ como el
intermedio PTG.3 deja sin cumplir la influencia de ese intermedio en
el próximo estado que exige AIS §381. Usar $F$ adicionalmente como
entrada al núcleo contravendría la regla de uso único de 90C §6.3(7).
La evaluación de RBG3(XOR) tiene un resultado NIST propio
\cite[§§375, 380--383]{ais}\cite[§§6.3(7), 6.4.1.2 y 6.4.2]{nist90c}.

\begin{center}
\small
\begin{tabularx}{\linewidth}{@{}lYY@{}}
\toprule
Operación & Consumo físico por operación en el perfil & Resultado contabilizado \\
\midrule
AES/SHA directo & Un $A$ de 128 bits. Inicio o resiembra aparte con 256 bits intermedios. & Un bloque externo de 128 bits. \\
CTR directo & Un $A$ de 128 bits. Material aparte de 384 bits full entropy para inicio o resiembra. & Un bloque externo de 128 bits. \\
CMAC en la rama física & Dos bloques Toeplitz de 128 bits. Justificar al menos 192 bits de entropía acreditada por evaluación. & Un tag completo de 128 bits, biyectivo en el segundo bloque dado el primero. \\
XOR con CTR & Un bloque físico full entropy de 128 bits por salida. Material de resiembra separado antes de continuar al alcanzar $2^{17}$ bits del DRBG. & 128 bits en el puerto XOR, con su propia ruta 90C. \\
RS con CTR & 384 bits full entropy por iteración. Si se producen con la variante CMAC descrita, seis bloques Toeplitz, 768 bits de entrada y al menos $3\cdot192=576$ bits acreditados. & 256 bits generados y 128 devueltos para esta petición. \\
\bottomrule
\end{tabularx}
\end{center}
Las cantidades acreditadas se sostienen con el presupuesto del puerto
90B. Las tasas raw se miden por separado. Cada coste de siembra/renovación es
adicional al de las muestras destinadas al selector y a los datos.
La regla de $2^{17}$ bits también se aplica en los supuestos de RBG2(P)
§5.3(6). El coste de seis bloques para la resiembra corresponde al CMAC
de dos bloques por tag y depende de ese diseño de recolector.

\paragraph{Correspondencia formal.}
El módulo \proofmodule{NISTProfiles} opera con palabras de 128 bits, claves de
256 y un cifrador con las dos leyes de inversa de una permutación.
\proofat{NISTProfiles.rsRecoversLow}{rsRecoversLow} y \proofat{NISTProfiles.rsProducesEveryBlock}{rsProducesEveryBlock} prueban ambas
composiciones de la inversa anterior. \proofat{NISTProfiles.cmacRecoversLast}{cmacRecoversLast} y
\proofat{NISTProfiles.cmacProducesEveryTag}{cmacProducesEveryTag} comprueban la coordenada final de un
mensaje de dos bloques completos, y \proofat{NISTProfiles.cmacDirectRoundtrip}{cmacDirectRoundtrip} y
\proofat{NISTProfiles.cmacDirectOnto}{cmacDirectOnto} la composición con el CTR directo. La función
\nolinkurl{rsStep} registra también el estado después de los dos bloques y
del Update posterior. El teorema \proofat{NISTProfiles.consecutiveDistinct}{consecutiveDistinct} verifica que los dos bloques difieren.\par
En el puerto XOR, \proofat{NISTProfiles.xorRecoversPhysical}{xorRecoversPhysical}
recupera $F$ dada la máscara. La invariancia del estado siguiente al variar $F$
se comprueba en \proofat{NISTProfiles.xorNextStateIndependent}{xorNextStateIndependent}.
Estas identidades usan la
permutación por clave. La acreditación de entropía física y la aprobación
de los mecanismos corresponden a la evaluación del producto.
\par
